\scanpage{21}\mspage{42}
on trouve une suit\deleted{\unsure{é}}e exacte canonique \handwritten{(où $\mathbb Z_l[1]=\varprojlim_{\nu}\deleted{\unsure{Z}}\mu_{l^\nu}$ est le faisceau $l$-adique de TATE) :}
\[
0\handwritten{\longrightarrow}t^{i-1}(X,l)\handwritten{\longrightarrow}H^i(X,\underline{\mathrm Z}_l\handwritten{[1]})\handwritten{\longrightarrow}T_l(H^i(X,\handwritten{\mu}_{l^\infty}))\handwritten{\longrightarrow}0\ , \tag*{(8.3)}
\]
où \deleted{$\unsure{t|H}^{i-1}(X,l)$} on a posé
\[
t^{i-1}(X,l)=H^{i-1}(X,\handwritten{\mu}_{\deleted{\unsure{l}}l^\infty})/H^{i-1}(X,\handwritten{\mu}_{l^\infty})^{\unsure{o|0}} \tag*{(8.4)}
\]
qui est donc un groupe fini de $l$-torsion, et s'identifie donc, par (8.3), au sous-\deleted{groupe}\inserted{module} de torsion d\deleted{\unsure{a|u}}\unsure{u|a} $\underline{\mathrm Z}_l$-module de type fini
\[
H^i(X,\underline{\mathrm Z}_l\handwritten{[1]})=\lim H^i(X,\deleted{\unsure{\underline{\mathrm Z}/l^\nu}}\handwritten{\mu_{l^\nu}})\ , \tag*{(8.5)}
\]
dont la partie libre s'identifie, à son tour, à $T_l(H^i(X,\handwritten{\mu}_{l^\infty}))$. On en \margin{\handwritten{!}}déduit également un isomorphisme
\[
\deleted{\unsure{H}}(H^i(X,\handwritten{\mu}_{l^\infty})^o\handwritten{\simeq} H^i(X,\underline{\mathrm Z}_l\handwritten{[1]})\otimes_{\underline{\mathrm Z}_l}\underline{\mathrm Q}_l\ /\ \mathrm{Im}\ H^i(X,\underline{\mathrm Z}_l\handwritten{[1]}), \tag*{(8.6)}
\]
qui montre que $H^i(X,\handwritten{\mu}_{l^\infty})$ et $H^i(X,\underline{\mathrm Z}_l\handwritten{[1]})$ se déterminent mutuellement mod groupes \deleted{\unsure{dont}} finis.

\handwritten{8.2.} Un passage à la limite analogue dans l\deleted{es}a suite\deleted{s} exacte\deleted{s} de Kummer (BR II 3.1) nous donne de même une suite exacte
\[
0\handwritten{\longrightarrow}\mathrm{NS}(X)\otimes_{\underline{\mathrm Z}}\underline{\mathrm Z}_l\handwritten{\longrightarrow}H^2(X,\underline{\mathrm Z}_l[1])\handwritten{\longrightarrow}T_l(\mathrm{Br}'(X))\handwritten{\longrightarrow}0\ , \tag*{(8.7)}
\]
en supposant maintenant que $\mathrm{Pic}(X)$ est extension d'un groupe de type fini $\mathrm{NS}(X)$ (le \uline{groupe de Néron-Sévéri}) par un groupe divisible, et que $H^2(X,\underline{\mathrm Z}/l\underline{\mathrm Z})$ est fini. \deleted{\unsure{(}}On trouve en passant le fait, classique dans la théorie transcendante, que \deleted{le}\inserted{la composante $l$-primaire du} sous-groupe de torsion de $\mathrm{NS}(X)$ est is\deleted{\unsure{o}}omorphe au sous-groupe de t\deleted{\unsure{ro}}orsion de $H^2(X,\underline{\mathrm Z}_l[1])$, en même temps qu'on précise du point de vue de la cohomologie $l$-adique la relation entre le groupe de Brauer (ou plus précisément, $T_l(\mathrm{Br}'(X))$) et la\textquotedblleft partie transcendante\textquotedblright\ du $H^2$ \handwritten{,} \deleted{\unsure{;}}comparer (BR II 3 passim ).

Lorsque $X$ est \inserted{une surface} propre, lisse et connexe sur un corps $k$ algébriquement clos, il y a lieu de considérer sur le terme médian de (8.7) la forme bilinéaire symétrique canonique, fournie par le cup-produit,\deleted{\unsure{à}} qui est à valeurs d\deleted{\unsure{q}}ans
% archive sheet number: 20
